% Options for packages loaded elsewhere
\PassOptionsToPackage{unicode}{hyperref}
\PassOptionsToPackage{hyphens}{url}
\PassOptionsToPackage{dvipsnames,svgnames,x11names}{xcolor}
\documentclass[
  10pt,
]{article}
\usepackage{xcolor}
\usepackage[margin=1cm,top=1cm,bottom=1cm,left=1cm,right=1cm,includeheadfoot]{geometry}
\usepackage{amsmath,amssymb}
\setcounter{secnumdepth}{3}
\usepackage{iftex}
\ifPDFTeX
  \usepackage[T1]{fontenc}
  \usepackage[utf8]{inputenc}
  \usepackage{textcomp} % provide euro and other symbols
\else % if luatex or xetex
  \usepackage{unicode-math} % this also loads fontspec
  \defaultfontfeatures{Scale=MatchLowercase}
  \defaultfontfeatures[\rmfamily]{Ligatures=TeX,Scale=1}
\fi
\usepackage{lmodern}
\ifPDFTeX\else
  % xetex/luatex font selection
  \setmainfont[]{DejaVu Serif}
  \setmonofont[]{DejaVu Sans Mono}
\fi
% Use upquote if available, for straight quotes in verbatim environments
\IfFileExists{upquote.sty}{\usepackage{upquote}}{}
\IfFileExists{microtype.sty}{% use microtype if available
  \usepackage[]{microtype}
  \UseMicrotypeSet[protrusion]{basicmath} % disable protrusion for tt fonts
}{}
\usepackage{setspace}
\makeatletter
\@ifundefined{KOMAClassName}{% if non-KOMA class
  \IfFileExists{parskip.sty}{%
    \usepackage{parskip}
  }{% else
    \setlength{\parindent}{0pt}
    \setlength{\parskip}{6pt plus 2pt minus 1pt}}
}{% if KOMA class
  \KOMAoptions{parskip=half}}
\makeatother
\usepackage{listings}
\newcommand{\passthrough}[1]{#1}
\lstset{defaultdialect=[5.3]Lua}
\lstset{defaultdialect=[x86masm]Assembler}
\usepackage{longtable,booktabs,array}
\usepackage{caption}
\captionsetup[table]{skip=6pt}
\newcounter{none} % for unnumbered tables
\usepackage{calc} % for calculating minipage widths
% Correct order of tables after \paragraph or \subparagraph
\usepackage{etoolbox}
\makeatletter
\patchcmd\longtable{\par}{\if@noskipsec\mbox{}\fi\par}{}{}
\makeatother
\setlength{\emergencystretch}{3em} % prevent overfull lines
\providecommand{\tightlist}{%
  \setlength{\itemsep}{0pt}\setlength{\parskip}{0pt}}
% Enable graphics inclusion and ensure figure numbering works
\usepackage{graphicx}
\renewcommand{\figurename}{Figure}

% Configure fonts for Unicode support with fallbacks
\usepackage{newunicodechar}
\newunicodechar{⁴}{\textsuperscript{4}}
\newunicodechar{₄}{\textsubscript{4}}

% Math notation macros
\newcommand{\norm}[1]{\lVert #1\rVert}
\newcommand{\abs}[1]{\lvert #1\rvert}

% Enhanced code block styling for better contrast and readability
\usepackage{fancyvrb}
\usepackage{xcolor}
\usepackage{listings}

% Define custom colors for code blocks
\definecolor{codebg}{RGB}{245, 245, 245}      % Light gray background
\definecolor{codeborder}{RGB}{200, 200, 200}  % Medium gray border
\definecolor{codefg}{RGB}{50, 50, 50}         % Dark gray text

% Configure Verbatim environment for inline code
\DefineVerbatimEnvironment{Verbatim}{Verbatim}{%
    fontsize=\small,
    frame=single,
    framerule=0.5pt,
    framesep=3pt,
    rulecolor=\color{codeborder},
    bgcolor=\color{codebg},
    fgcolor=\color{codefg}
}

% Configure code block styling
\DefineVerbatimEnvironment{Highlighting}{Verbatim}{%
    fontsize=\footnotesize,
    frame=single,
    framerule=0.5pt,
    framesep=5pt,
    rulecolor=\color{codeborder},
    bgcolor=\color{codebg},
    fgcolor=\color{codefg}
}

% Style inline code with \texttt
\renewcommand{\texttt}[1]{%
    \colorbox{codebg}{\color{codefg}\ttfamily #1}%
}

% Configure listings package for code blocks
\lstset{
    backgroundcolor=\color{codebg},
    basicstyle=\footnotesize\ttfamily\color{codefg},
    breakatwhitespace=false,
    breaklines=true,
    captionpos=b,
    commentstyle=\color{codefg},
    deletekeywords={...},
    escapeinside={\%*}{*)},
    extendedchars=true,
    frame=single,
    framerule=0.5pt,
    framesep=5pt,
    keepspaces=true,
    keywordstyle=\color{codefg},
    language=Python,
    morekeywords={*,...},
    numbers=left,
    numbersep=5pt,
    numberstyle=\tiny\color{codefg},
    rulecolor=\color{codeborder},
    showspaces=false,
    showstringspaces=false,
    showtabs=false,
    stepnumber=1,
    stringstyle=\color{codefg},
    tabsize=2,
    title=\lstname
}

% Override any Pandoc default lstset configurations
\AtBeginDocument{
    \lstset{
        backgroundcolor=\color{codebg},
        basicstyle=\footnotesize\ttfamily\color{codefg},
        frame=single,
        framerule=0.5pt,
        framesep=5pt,
        rulecolor=\color{codeborder},
        numbers=left,
        numbersep=5pt,
        numberstyle=\tiny\color{codefg}
    }
}

% Configure hyperref colors consistently
\AtBeginDocument{
% Override pandoc's hidelinks setting with consistent options
\hypersetup{
    colorlinks=true,
    allcolors=red,
    linkcolor=red,
    urlcolor=red,
    citecolor=red,
    filecolor=red,
    menucolor=red,
    linktoc=all
}
}

% Simple page break support for document structure
\usepackage{bookmark}
\IfFileExists{xurl.sty}{\usepackage{xurl}}{} % add URL line breaks if available
\urlstyle{same}
% fallback for those not using the hyperref driver hyperxmp:
\makeatletter
\@ifundefined{xmpquote}{\newcommand{\xmpquote}[1]{#1}}{}
\makeatother
\hypersetup{
  colorlinks=true,
  linkcolor={red},
  filecolor={red},
  citecolor={red},
  urlcolor={red},
  pdfcreator={LaTeX via pandoc}}

\title{Lattice Tooling: Omnidirectional Numbering and Nearest-Site Search}
\author{Daniel Ari Friedman\\ ORCID: 0000-0001-6232-9096\\ Email: daniel@activeinference.institute\\ DOI: 10.5281/zenodo.16887791}
\date{October 07, 2026}

\begin{document}
\maketitle

{
\hypersetup{linkcolor=black}
\setcounter{tocdepth}{3}
\tableofcontents
}
\setstretch{1.0}
\section{Lattice Tooling: Omnidirectional Numbering and Nearest-Site
Search}\label{lattice-tooling-omnidirectional-numbering-and-nearest-site-search}

\subsection{Overview}\label{overview}

Sections \href{02_4d_namespaces.md}{2} and
\href{03_quadray_methods.md}{3} develop the Quadray/IVM framework
analytically; this section documents the computational tooling that
operates directly on the close-packed lattice. Two modules extend the
analytical core without touching it:

\begin{itemize}
\tightlist
\item
  \passthrough{\lstinline!omni\_numbering!} --- omnidirectional
  close-packing numbering: vectorized enumeration of the IVM shell
  sequence, cumulative counts, and bidirectional site/index mappings.
\item
  \passthrough{\lstinline!lattice\_search!} --- fast
  nearest-lattice-point queries:
  \passthrough{\lstinline!nearest(site, R, k)!} and
  \passthrough{\lstinline!within\_radius(site, R)!} built on a
  precomputed ball index, entirely in NumPy (no scipy).
\end{itemize}

Both modules operate on canonical quadray integer 4-tuples (non-negative
components, at least one zero --- the projective normalization of
\href{03_quadray_methods.md}{Section 3}), and both rely on the same
verified shell invariants.

\subsection{Omnidirectional Close Packing and Frequency
Shells}\label{omnidirectional-close-packing-and-frequency-shells}

In the isotropic vector matrix (IVM), equal spheres pack
omnidirectionally --- closest packing. Starting from a central sphere,
the packing builds up in consecutive \textbf{frequency shells}: the
shell of frequency \(k \in \mathbb{Z}_{\ge 0}\) sits \(k\) radius-length
increments outward of the center, the same four-dimensional buildup
language synergetics uses for the growing vector equilibrium. Writing
\(N_k\) for the population of shell \(k\), the lone center is the whole
of shell 0 (\(N_0 = 1\)), and every shell with \(k \ge 1\) carries
exactly

\begin{equation}\label{eq:lattice-shell-population}
N_k \;=\; 10\,k^{2} \;+\; 2
\end{equation}

sphere centers. The first shell --- the 12 permutations of
\((2, 1, 1, 0)\) --- is the cuboctahedron (vector equilibrium) of the
``twelve around one'' motif; each subsequent shell is the next layer of
the omnidirectional buildup. The cumulative count through frequency
\(k\), written \(C_k\), is the centered-cuboctahedral closed form

\begin{equation}\label{eq:lattice-cumulative}
C_k \;=\; 1 \;+\; \sum_{j=1}^{k} \left(10\,j^{2} + 2\right) \;=\; 1 + 2k + \frac{10\,k\,(k+1)\,(2k+1)}{6}.
\end{equation}

The functions \passthrough{\lstinline!shell\_count(k)!} and
\passthrough{\lstinline!cumulative\_count(k)!} evaluate Eqs.
\eqref{eq:lattice-shell-population} and \eqref{eq:lattice-cumulative}
(scalar or array input), and the test suite cross-validates them against
direct summation.

\subsubsection{Reachability is enumerative, not
congruence-based}\label{reachability-is-enumerative-not-congruence-based}

A normalized integer quadray \(q\) with \(\sum_i q_i \equiv 0 \pmod 4\)
is a candidate IVM point, but not every candidate is reachable from the
origin by the 12 neighbor moves. The module therefore builds the site
set \textbf{layer by layer}: each new shell is produced at once by
adding all 12 moves to the previous frontier, re-normalizing rows, and
removing sites already seen (packed-key membership tests). The result is
cached per depth: \passthrough{\lstinline!sites\_through\_shell(k)!}
returns the \passthrough{\lstinline!(C\_k, 4)!} integer array in
canonical order --- the center first, then each shell in lexicographic
\passthrough{\lstinline!(a, b, c, d)!} order --- and
\passthrough{\lstinline!generate\_shell(k)!} returns one shell. An
independent breadth-first reference over
\passthrough{\lstinline!itertools.permutations!} confirms exact
agreement through frequency 8, including the counts of Eq.
\eqref{eq:lattice-shell-population} on every shell.

\subsubsection{Site/index mapping}\label{siteindex-mapping}

\passthrough{\lstinline!site\_index(site, max\_shell)!} returns the
position of a site in that canonical enumeration (or
\passthrough{\lstinline!-1!} if absent within the depth), and
\passthrough{\lstinline!site\_at\_index(index, max\_shell)!} is its
inverse. Lookups accept any projective representative --- the input is
translated by \passthrough{\lstinline!-(k, k, k, k)!} before matching
--- so \passthrough{\lstinline!(2, 1, 1, 0)!} and
\passthrough{\lstinline!(3, 2, 2, 1)!} map to the same index, while
\passthrough{\lstinline!(0, 1, 1, 2)!}, a different permutation and
hence a different sphere, maps elsewhere. Component overflow is guarded:
inputs whose normalized components reach the 16-bit packed-key width
(\(2^{16}\), the per-component field of the int64 keys) return
\passthrough{\lstinline!-1!} rather than wrapping.

Executable check of the bookends:

\begin{lstlisting}[language=Python]
from quadmath.lattice.omni_numbering import (
    shell_count, cumulative_count, sites_through_shell,
    site_index, site_at_index,
)

assert shell_count(1) == 12 and shell_count(2) == 42      # 10k^2 + 2
assert cumulative_count(2) == 55                          # 1 + 12 + 42
assert site_index(site_at_index(12, 3), 3) == 12          # roundtrip
assert site_index((2, 0, 0, 0), 6) == -1                  # unreachable site
\end{lstlisting}

\subsection{Fast Nearest-Site Queries}\label{fast-nearest-site-queries}

\subsubsection{Exact lattice distances}\label{exact-lattice-distances}

Under \passthrough{\lstinline!quadray.DEFAULT\_EMBEDDING!} the four
embedding columns \(C_j\) satisfy \(C_i \cdot C_j = 4\,\delta_{ij} - 1\)
at unit scale, so the Euclidean squared distance between integer
4-vectors \(p, q\) is the exact integer

\begin{equation}\label{eq:lattice-distance-identity}
d^{2}(p, q) \;=\; 4 \sum_{j} \delta_j^{2} \;-\; \Bigl(\sum_{j} \delta_j\Bigr)^{2},
\qquad \delta = p - q,
\end{equation}

implemented in \passthrough{\lstinline!squared\_distance!} and
cross-validated in the tests against direct embedding coordinates.
Because the identity is exact on integers, ranking and filtering never
suffer floating-point boundary errors: a site is inside the query radius
\(R > 0\) iff \(d^{2} \le R^{2}\) evaluated exactly (in float64 the same
algebra is exact for the magnitudes involved here).

\subsubsection{Truncation bound}\label{truncation-bound}

Every site \(s\) on frequency shell \(g\) satisfies
\(d^{2}(\mathbf{0}, s) \ge 8\,g\) --- verified by exhaustive enumeration
through frequency 8 in the test suite --- and the shell maximum is
exactly \(8\,g^{2}\). Combining this invariant with the triangle
inequality gives the depth at which a query can stop:

\begin{equation}\label{eq:lattice-truncation-bound}
g \;\leq\; \frac{\left(\lVert c \rVert + R\right)^{2}}{8},
\end{equation}

for a query center \(c\) and radius \(R > 0\), where \(\lVert c \rVert\)
is the Euclidean norm of the embedded center: any site within \(R\) of
\(c\) must live on a shell at most this deep.
\passthrough{\lstinline!within\_radius(site, R)!} therefore enumerates
exactly through \(\lfloor (\lVert c \rVert + R)^{2} / 8 \rfloor + 1\)
shells, filters with Eq. \eqref{eq:lattice-distance-identity}, and
returns all hits sorted by squared distance with lexicographic
\passthrough{\lstinline!(a, b, c, d)!} tie-breaking. Requests whose
required depth exceeds the precomputed index depth
(\passthrough{\lstinline!MAX\_SHELL = 32!}, about \(1.1 \times 10^{5}\)
sites --- \(C_{32} = 114465\) by Eq. \eqref{eq:lattice-cumulative})
raise \passthrough{\lstinline!ValueError!} rather than silently
truncating.

\subsubsection{\texorpdfstring{Shell-sweep
\texttt{nearest}}{Shell-sweep nearest}}\label{shell-sweep-nearest}

\passthrough{\lstinline!nearest(site, R, k)!} sweeps shells outward,
ranking each shell's squared distances with
\passthrough{\lstinline!numpy.argsort!}; once
\passthrough{\lstinline!k!} candidates are in hand, the invariant behind
Eq. \eqref{eq:lattice-truncation-bound} supplies an early-stop test ---
no unvisited shell can contain a closer site: a site at depth \(g + 1\)
satisfies \(d^{2}(\mathbf{0}, s) \ge 8\,(g+1)\), so the sweep stops once
\(8\,(g+1) > (\sqrt{d^{2}_{(k)}} + \lVert c \rVert)^{2}\), with
\(d^{2}_{(k)}\) the current k-th best squared distance. The sweep never
under-enumerates (the bound is conservative; ties at the boundary are
kept), and answers always match the brute-force reference over the same
region, as the tests assert on fixed-seed random queries.

\begin{lstlisting}[language=Python]
import numpy as np
from quadmath.lattice.lattice_search import nearest, within_radius

sites, d2 = nearest((0.4, -0.3, 0.9, 0.1), R=2.0, k=5)   # 5 nearest centers
ball, ball_d2 = within_radius((2, 1, 1, 0), R=4.0)       # everything within 4
assert np.all(np.diff(ball_d2) >= 0)                      # ascending distances
\end{lstlisting}

\subsection{API Summary}\label{api-summary}

{\def\LTcaptype{none} % do not increment counter
\begin{longtable}[]{@{}
  >{\raggedright\arraybackslash}p{(\linewidth - 2\tabcolsep) * \real{0.5000}}
  >{\raggedright\arraybackslash}p{(\linewidth - 2\tabcolsep) * \real{0.5000}}@{}}
\toprule\noalign{}
\begin{minipage}[b]{\linewidth}\raggedright
Function
\end{minipage} & \begin{minipage}[b]{\linewidth}\raggedright
Purpose
\end{minipage} \\
\midrule\noalign{}
\endhead
\bottomrule\noalign{}
\endlastfoot
\passthrough{\lstinline!omni\_numbering.shell\_count(k)!} & Shell
population, Eq. \eqref{eq:lattice-shell-population} \\
\passthrough{\lstinline!omni\_numbering.cumulative\_count(k)!} &
Cumulative count through shell \passthrough{\lstinline!k!}, Eq.
\eqref{eq:lattice-cumulative} \\
\passthrough{\lstinline!omni\_numbering.generate\_shell(k)!} & Sites of
one frequency shell (lexicographic) \\
\passthrough{\lstinline!omni\_numbering.sites\_through\_shell(k)!} &
Whole enumeration \passthrough{\lstinline!(C\_k, 4)!} in canonical
order \\
\passthrough{\lstinline!omni\_numbering.site\_index!} /
\passthrough{\lstinline!site\_at\_index!} & Bidirectional site/index
mapping with projective normalization \\
\passthrough{\lstinline!lattice\_search.squared\_distance!} & Exact
squared distances via Eq. \eqref{eq:lattice-distance-identity} \\
\passthrough{\lstinline!lattice\_search.within\_radius(site, R)!} & All
sites within radius, exact and sorted \\
\passthrough{\lstinline!lattice\_search.nearest(site, R, k)!} &
\passthrough{\lstinline!k!} nearest sites within radius, shell-sweep
with early stop \\
\end{longtable}
}

\subsection{Verification}\label{verification}

Both modules ship with 100\% test coverage (branch coverage included):
the enumeration is checked against an independent breadth-first
reference through frequency 8, distance identities against direct
embedding computations, and query results against brute-force filtering
on fixed-seed random centers. Run:

\begin{lstlisting}[language=bash]
PYTEST_DISABLE_PLUGIN_AUTOLOAD=1 uv run coverage run -m pytest -q
uv run coverage report
\end{lstlisting}


\end{document}
